\documentclass{article}
\usepackage[russian]{babel}
\usepackage[T2A]{fontenc}
\usepackage[utf8]{inputenc}
\usepackage{amsmath}

\title{О сходимости рядов Фурье}
\author{И.~Петров}
\date{2026}

\begin{document}
\maketitle

\section{Введение}

Теория рядов Фурье занимает центральное место в математическом анализе.
Пусть $f \colon [-\pi, \pi] \to \mathbb{R}$ --- интегрируемая функция.
Ряд Фурье функции~$f$ определяется как
\[
  f(x) \sim \frac{a_0}{2} + \sum_{n=1}^{\infty} \left(a_n \cos nx + b_n \sin nx\right),
\]
где коэффициенты вычисляются по формулам
\[
  a_n = \frac{1}{\pi}\int_{-\pi}^{\pi} f(x)\cos nx\,dx, \quad
  b_n = \frac{1}{\pi}\int_{-\pi}^{\pi} f(x)\sin nx\,dx.
\]

\section{Основной результат}

\begin{theorem}[Карлесон]
Если $f \in L^2[-\pi,\pi]$, то ряд Фурье функции~$f$ сходится почти всюду.
\end{theorem}

Доказательство этой теоремы является одним из наиболее сложных результатов
гармонического анализа двадцатого века.

\section{Заключение}

Проблема сходимости рядов Фурье продолжает оставаться активной областью
исследований в современной математике.

\end{document}
